
% 来源：EPUB part091.xhtml
\section{好的企业文化就是做对的事情}

所谓企业文化，讲的就是什么是对的事情（或者说哪些是不对的事情），以及如何把对的事情做对。苹果的企业文化堪称好的企业文化的经典，多看看苹果的发布会或许能有所悟。（2015-03-22）

% 来源：EPUB part092.xhtml
\subsection{文化好的企业活得长些}

企业文化作为过滤器非常有威力，为我避免了很多错误。怎么选对的公司是能力问题，不选错的公司是是非问题。（2019-06-06）

好的企业文化可以为好的商业模式保驾护航，甚至可以帮助你发现好的商业模式，但我没想过权重的问题。这是个过滤器，没有权重的概念。（2019-05-23）

我不会只是基于企业文化买公司，但会把有好的企业文化作为想买的前提之一。（2015-05-25）

\BookQuestion{网友：我会把有“不好”的企业文化的公司作为不买的全部理由。}

确实如此，不好的公司文化时间长了一定会伤害到企业自身，而且不好的企业文化一旦建立就很难去掉了。（2015-05-28）

\BookQuestion{网友：您最看重企业文化的哪些方面？}

我看重那些让企业更健康更长久的因。（素2019-04-09）

\BookQuestion{网友：好的生意模式需要有好的企业文化，只有好的企业文化不一定有好的生意模式，是这样吗？}

对，特别好的生意模式也是可遇不可求的，而且机会一旦错过就没了。不过，没有好的生意模式也是有可能成为一家不错的企业的。（2011-12-17）

\BookQuestion{网友：苹果算不算傻瓜也可以经营好的公司？}

其实我并不知道彼得·林奇的定义到底是什么意思，因为现实里没有什么企业是傻瓜摧毁不了的，关键是看给这个傻瓜多长时间。我猜，这里傻瓜都能经营的企业指的是这个企业既有着强大的商业模式，能经得起傻瓜折腾“一小会儿”，同时有好的企业文化能够很快纠偏。就算好的企业经营者也会有犯傻的时候，但好的企业文化会让企业在比较短的时间里改正。比如当年谷歌买了摩托罗拉，结果很快就发现不合适，然后一年多就改了。（2024-05-17）

所谓好的企业文化大概指的就是“利润之上的追求”，而“利润至上”绝不可能成为好的企业文化。好的企业文化未必能形成好的生意模式。个人觉得好的生意模式是一定要有好的企业文化做支撑的，所以是必要条件。但好的企业文化不是拥有好的生意模式的充分条件。（2013-02-12）

企业文化确实很重要，这点上我比巴菲特要更强调。（2015-04-29）

\BookQuestion{网友：巴菲特说对的管理层，大道强调企业文化，这两者有何区别？}

现在巴菲特也说的是企业文化了。（2019-03-18）

可能是因为我想早退休，所以总说企业文化，巴菲特一直是CEO，所以说的是管理层？现在他应该在考虑退休以后的事情了，所以也经常说企业文化了？我猜着玩的，实际上我们讲的是一个东西，没本质差别。（2019-06-23）

\BookQuestion{网友：企业文化和规矩的关系是什么？}

简单讲就是规矩管不着的地方文化管。（2016-07-05）

\BookQuestion{网友：企业文化这个词有点虚无缥缈，我可不可以理解成马云说的使命感、价值观、愿景、组织、人才、KPI这样的东西？要通过什么维度来判断企业文化好呢？}

企业文化其实非常实在，建议你看看苹果的发布会，体会一下不虚无缥缈的企业文化。建议你从现在往前看10年的发布会，看完你会有体会的。（2020-12-04）

所谓的企业文化好并不是百战百胜的武器，他只是能让企业少犯原则性错误而已。或者说，在同一个行业里，平均而言企业文化好的企业活得长些而已。（2011-02-11）

企业文化不好的公司长期必然日子没那么好过，好企业偶然也会有难过的日子。（2023-05-14）

\BookQuestion{网友：您怎么看待战略和企业文化的关系？}

有好的企业文化的公司战略出错的概率低。（2011-08-25）

有好的企业文化的公司往往应变能力要强很多。（2010-04-05）

\BookQuestion{网友：维持好生意，主要是靠好管理层建立的好企业文化吗？}

没有好的企业文化很难维持，但有好的企业文化也未必就一定能维持住的。因为也可能会有战略错误导致企业完蛋的，尤其是小企业。（2013-04-20）

\BookQuestion{网友：茅台生意属性很好，但企业文化未必好，如果二者不能兼得该如何取舍？}

如果茅台的企业文化不够好的话，茅台走不到今天的。这不等于说茅台没一点问题。企业文化不好的公司，好的商业模式早晚是维持不了的。（2019-04-25）

\BookQuestion{网友：您认同企业文化跟您做企业的经验有关系。可以因为一家公司是百年老店而判断它未来也一定是好的吗？}

百年老店不意味着它就一定能活到101年，这就像不能用市盈率去预测明年的利润一样。雷曼好像就有150多年的历史，但只要你注意到雷曼后期的企业文化都变成什么样了的话，你就不会对它的结局感到惊奇。不过，有好的企业文化的企业犯原则性大错误的机会低，和用户比较近，表现好的机会大。《基业长青》里会讲得比较好些。（2010-04-04）

一般来讲，公司的企业文化是由三个部分组成的：使命、愿景以及核心价值观。使命指的是企业存在的意义；愿景是企业内大家的共同远景。（2010-04-03）

我觉得企业文化里最重要的应该是核心价值观的问题，不符合核心价值观的东西不做最重要。（2020-10-20）

\BookQuestion{网友：国内有些企业似乎很看重且一直在强调自己的市场占有率，这样的公司段总会不会直接忽略掉？是不是只专注于消费者的公司会好一些？}

讲本身并不是问题，作为一个大的愿景讲出来没什么不好，尤其是企业初期，有一个大的愿景非常重要。差别在于用什么办法去达到。用消费者导向去达到愿景目标是本手。（2010-10-24）

愿景的意思是大家的共同远景，是可实现的。（2011-10-02）

没有愿景的公司容易陷入利润导向，最后容易被眼前利益诱惑而犯大错，从而导致公司寿命缩短。

\BookQuestion{网友：很多公司都有愿景，但好像都是写在字面上的东西。我观察公司愿景是真是假，主要看领导有没有真的把它当回事。您经常提到步步高的愿景，所以我相信它是真的。}

愿景的意思是大家的愿景（远景），光写在纸上是没用的。我们公司提的是更健康更长久。也就是说不健康不长久的事我们不应该做（哪怕眼前有钱赚也不应该做）。（2010-11-10）

\BookQuestion{网友：有了愿景不一定会成功，但没有愿景，一般不会成功。}

那要看如何定义成功。没有愿景的公司很难走得很远，但也有些小的家族企业可以靠别的东西维持走很远的。（2010-11-12）

\BookQuestion{网友：能举个例子吗？}

那你肯定不知道耐克的故事。值得读一下。耐克还是一个很小的公司时，曾经有一个远景，“打败阿迪达斯”，真是太可笑了（阿迪达斯当时是这么想的）。（2010-05-30）

\BookQuestion{网友：管理的最高境界是没有管理，制度的最高境界是没有制度！有统一的价值观和统一的目标，员工能自我约束，制度健全与否没有太大的重要性？}

没有管理公司很快就会完蛋的。所谓好的企业文化是能管到制度管不过来的东西，不是个神秘的咒语。（2018-11-07）

制度是强制性的，文化则不完全是，所以建立好的企业文化非常难，破坏起来非常容易。

\BookQuestion{网友：破坏起来非常容易如何理解？}

比如建立信誉很难，破坏很容易。（2018-11-08）

\BookQuestion{网友：您觉得企业文化可以改变或者能改变吗（由不好转好）？}

想想企业文化怎么来的就知道可不可以改变了。（2013-12-10）

\BookQuestion{网友：有些公司经历一些教训后，企业文化有可能慢慢改良吗？}

文化有点像习惯，也有点像基因。（2013-05-07）

\BookQuestion{网友：体制因素和其带来的对于企业经营及企业文化的影响，是否是更多投资于国内企业的投资者应该多注意的呢？}

可能我在体制内的企业待过，所以比较敏感，知道有些什么区别。这种区别不是随时都在，也不是哪里都能显现出来，但往往关键的时候就会发挥一些看不见的作用。（2010-08-14）

% 来源：EPUB part093.xhtml
\subsection{最重要的是什么不可以做}

\BookQuestion{网友：能用一句话总结一下什么是好的企业文化吗？}

做对的事情！（2019-05-29）

\BookQuestion{网友：好的企业文化应该是怎样的？}

也许你自己可以先看看哪些属于做对的事情，哪些属于把事情做对。另外，文化里的东西都是用来约束自己的，千万别用作挑剔别人的工具。企业文化的东西是需要整体认同的，不认同的人会离开或被赶走。没有任何一个所谓的好的企业文化可以适合所有的企业或组织。所谓合适性就是指的员工对企业文化是否适应，而合格性则是指把事情做对的能力。（2016-06-05）

\BookQuestion{网友：不太理解“千万别用来挑剔别人”。}

我早年说过，本分不是“照妖镜”，主要是用来检视自己。但企业同时有自己的不为清单，那个一般是强制性的。（2020-12-03）

\BookQuestion{网友：如何让字面上的企业文化落实到所有员工的言行中呢？}

最重要的是什么不可以做。你可以多看看韦尔奇的书。阿里巴巴在企业文化建立方面的水平也还不错。（2010-04-23）

韦尔奇的《赢》里讲了不少GE是怎么做的。阿里巴巴这方面做得非常好。这是一项几十年或更长的工作，应该和企业的寿命一样长，是一项重要但往往不那么紧迫、常常被人忽略的工作。（2010-04-04）

\BookQuestion{网友：什么是好的企业文化，如何建立好的企业文化，大道能推荐一本书吗？}

我觉得看书可能不如看企业来得明白。书可以看《基业长青》，企业可以看苹果这些年的产品发布会（每年9月份的那个就好），或者多看几遍巴菲特在大学的问答交流。（2019-06-01）

\BookQuestion{网友：您关于企业经营的知识是如何获取的？}

要搞清这个问题，你要回到什么是企业文化就明白了。我没有一套自己的关于企业经营的知识（我认为叫理解可能更合适些）。所有我对企业的理解，这本书或那本书里都有，但没有体会的人一般不会去体会而已。重要的是要有基本的原则：做对的事情，把事情做对。

在如何把事情做对上犯的错往往是不可避免且都还是有机会纠正的。（2011-01-25）

\BookQuestion{网友：段总在营销本质、企业文化方面的诠释是最简单通透的，一句话或一个例子，就让俺开悟了。}

哈，我们确实也有自己的体会，悟出来不容易，你能一下子体会到那绝对是你的造化，多数人做不到的。（2018-09-12）

好企业也可能会犯错的。我一般看到企业犯错时，主要判断这个错是做了错的事情，还是在把事情做对的过程当中犯的错。任何人在做对的事情的前提下，为了把事情做对的过程当中都是会犯错的，这种错误也可能给企业造成很大的短期损失，但长期而言却可能微不足道。而另一种错误的损失却可以是致命的。（2011-01-03）

我非常喜欢谷歌的“不作恶”（Don't be evil）。（2012-04-14）

\BookQuestion{网友：我在创业，不知道应该建立一个什么样的企业文化，您能不能给一些建议？}

企业文化必须你自己去创立，没人可以帮忙的。企业文化大致由使命、愿景和核心价值观组成，这都是你自己的东西。使命就是你的公司为什么存在，愿景就是大家共同的远景，核心价值观就是是非观，或者叫什么是对的事情，什么是不对的事情。（2013-06-06）

\BookQuestion{网友：步步高的企业文化是您创办的时候就已经想明白，还是创办之后慢慢形成的？}

有很多东西在创办之前就有了，但慢慢补充和完善或修改了一些。基本核心价值观是不会变的。（2013-03-26）

\BookQuestion{网友：成功企业往往会打上开创者的气质。}

那是因为企业文化一般都是由开创者建立的，当然，后期的完善也很重要。

\BookQuestion{网友：开创的时候企业文化怎么建立？我个人觉得企业文化起初的建立不需要投入多少精力吧？老板如何对待同事、员工、客户及其他合作伙伴，体现出来的这些作风就是企业文化的雏形吧？企业搞大了再整理成易理解的东西，首先传达给干部，干部传达战士。总之应该是一砖一瓦的事情。}

同意。（2010-04-04）

\BookQuestion{网友：企业文化就是老板文化，这个理解对不对？}

一般而言，创始人对企业文化的形成起的作用都很大。（2010-05-10）

\BookQuestion{网友：企业文化如何才能改变？}

大概办法就是年年讲月月讲天天讲，而且光讲是肯定没有用的，但不讲是绝对不会变的。（2023-12-12）

% 来源：EPUB part094.xhtml
\subsection{听其言观其行}

\BookQuestion{网友：如何判断企业文化好与坏，标准或者共性是什么？}

这个问题很大。笼统地讲，就是企业行事是以利益还是以是非为标准。如果凡事是以利益为准绳的，我就不太喜欢。

其实我也没有一个所谓的格式或公式去判别一个公司的企业文化，不过，我经常会用拟人化的角度去想一家公司。我不想打交道的人我也不想投资他们的公司。（2020-10-11）

很难判断一个企业的文化是否好，但看出不好的企业文化要容易得多。（2010-07-10）

基本上就是听其言观其行。看这家公司过去都说过啥，都怎么做的。比如，当年我对苹果开始有兴趣后，几乎看了苹果所有的发布会以及能找到的库克和乔布斯讲过的东西，也用了很多苹果的产品。（2020-10-19）

我关注的公司和人都很少，可以花足够时间在有限的人和公司上，这样准确率自然就高很多了。（2020-12-07）

我如果不信任公司的文化和管理层了，我会决定离开的。对我来说，要搞懂公司经营的每个细节来决定投反对票还是赞成票是件很费时的事情，比搞懂商业模式或者企业文化可能还要难。如果整体上我看不懂公司的某些重要行为，我多数会选择离开，虽然这样也会犯错，但比老是很膈应要舒服很多。（2023-06-14）

\BookQuestion{网友：我看企业文化的方法就是经常看领导人访谈。}

看访谈确实是个很重要的方法，结合过去他们都说过啥，是否都做到了，也可以看出来他们说的时候是真话还是套话。所以，看一个访谈是不够的，光看访谈也是不够的。（2010-11-09）

\BookQuestion{网友：怎么判断公司的企业文化和价值观是每个人都认定的，还是贴在办公室里的标语根本没人信？}

呵呵，你看企业的领导们信不信就可以了。如果最高领导人都不信，那肯定不行。（2010-04-04）

\BookQuestion{网友：如何观察公司这些写在纸面上的是不是真的愿景？}

看他的历史、说的和做的。（2010-11-12）

\BookQuestion{网友：绝大部分公司的企业文化，网站上的那短短百十个字都是差不多的。}

其实即便是字面也是差很多的。（2013-03-28）

产品文化是企业文化的一部分但不是全部。一般来讲能有挂在墙上的也比哪里都没有的强，但如果不是真的则强得有限。（2012-03-07）

\BookQuestion{网友：在步步高企业文化中股东居然放最后？如果真是按重要程度排序，这顺序不就是“消费者第一、员工第二、商业伙伴第三、股东第四”吗？}

前面三条是要做的事情，最后一条是我们的观点，就是说，我们认为做到了前面三条，股东回报应该不错。二十多年前我们写下这一条时就是这么认为的，迄今依然如此，结果似乎也正是如此。

我们不排序是因为我们认为也许不该排序。（2019-03-10）

\BookQuestion{网友：步步高企业文化的使命中有一条“对股东：使其投入的股本有高于社会平均报酬的合理回报”，“高于社会平均报酬”该如何理解？}

其实就是我们有信心比平均回报好的意思。

我觉得我们确实做到了哈，而且做得还不错！（2019-09-26）

\BookQuestion{网友：马云讲过，客户第一、员工第二、股东第三。您怎么看？}

反正倒过来摆是肯定不对的。（2010-05-20）

我可以接受他的意思。我们公司认为这三者关系是平等的。大家是从不同的角度去说这个问题。（2010-11-13）

\BookQuestion{网友：我认为马云不是投资人的好朋友。不尊重股东的人，也得不到股东的尊重！}

我倒是见到许多许多“股东第一”的公司，由于不停地为了取悦华尔街而采取短期行为，最后业绩差导致股东利益都没有了。马云这种实际上是最尊重股东利益的那一类了。我个人投资时比较喜欢这类“股东第三”的公司，最害怕那种“股东第一”的公司。（2012-01-02）

\BookQuestion{网友：马云说的客户第一、员工第二、股东第三，是不是有点“利润之上追求”这个意思。}

正是！（2012-07-24）

\BookQuestion{网友：什么样的企业文化是维护股东利益的企业文化？我一直有点模糊。}

维护股东利益指的是维护股东的长远利益，这里的长远指的是企业的整个生命周期，因为所谓的股东利益实际上就是企业未来净现金流的折现。（2013-03-26）

\BookQuestion{网友：难道不包括短期利益吗？}

长期利益中，本来就包括不损害长期利益的短期利益。（2013-03-28）

\BookQuestion{网友：阿里巴巴对于大股东的态度和处理与大股东关系的方法，让我感觉多少有些不确定性。}

我没觉得阿里巴巴集团对大股东的态度有什么问题啊。我倒是觉得他们的态度是上市公司里最好的。因为大股东真正需要的是公司能好好地专注在公司的发展上，而且大股东其实只需要这一条就够了。（2010-05-25）

\BookQuestion{网友：参照巴菲特的8项投资标准，不以所有的股东利益而是以大股东利益为先的企业，我们能指望什么呢？}

什么叫以大股东利益优先？股东利益只能通过投票决定。还有就是如果你不喜欢，觉得公司不以你的利益优先的话，你是应该卖的，不然你就是觉得他们会以股东利益优先，否则就矛盾了。有些大股东利用控股来为自己牟利，这种股票最好远离。

另外，股东利益优先和股东第三并不矛盾，是不同角度。（2010-05-28）

我反对任何公司（包括我们自己公司）给股东任何优惠，和茅台无关。作为比你们大很多的茅台股东，我很高兴看到茅台取消这么荒唐的东西，觉得茅台在这点上进步了。（2022-06-11）

\BookQuestion{网友：段先生这样说是直指大股东吧！大股东拿货是969元的出厂价，而小股东买的是零售价1499元，不占公司的便宜，而且还给公司增加利润。}

你是说大股东能赚便宜为什么我不能赚？我说的是谁都不应该赚公司的便宜，包括股东，当然也包括大股东。（2022-07-14）

\BookQuestion{网友：您如何看同股不同权？}

我们是同股同权的，我认为同股不同权不公平，但没仔细想过为什么会有这个东西。可能是有些公司怕被资本大鳄偷袭而产生的一个办法？感觉长期来说这个东西对公司是有害的。（2023-05-11）

\BookQuestion{网友：你们公司的人要成为公司股东得工作多少年才有资格？成为股东后，如果发现这个人并不适合再合作下去，你们是怎么做的？}

我觉得你对股东的定义是不是有什么误解？公司里一股的股东和一亿股的股东都叫股东，他们每股的权利是完全一样的。成为股东非常容易，只要有一股就是股东。公司和股东之间没有合作关系，但股东可以通过投票来决定董事会或者董事的人选以及公司策略等等……（2023-06-14）

\BookQuestion{网友：《基业长青》里面提到过，企业应该平衡员工、股东、消费者和合作伙伴的利益。最近发现有些企业比较压榨合作伙伴，但是也在比较长（10多年）的时间里面做得不错。这该怎样理解呢？}

第一，《基业长青》说的不是充要条件；第二，有时候10多年还不够长，我们邻居都三代了还混着呢。（2014-04-14）

% 来源：EPUB part095.xhtml
\subsection{案例}

\BookMinorHeading{格力电器}

\BookQuestion{网友：您认为格力公司如何？听说格力是先收款再发货的。}

格力是好公司。

呵呵，难道还有别的办法可以活下来的？（2010-03-30）

\BookQuestion{网友：有个人问董明珠，为什么最近格力的股票跌跌不休，她说我只对我的投资者负责，不对投机者负责。她给这几者的排序是：1. 消费者，2. 员工，3. 社会责任，4. 股东。}

董明珠对企业的理解非常好。（2010-05-27）

\BookQuestion{网友：决定买格力不是美的是因为一个标语。我去格力的时候，有一个厂房的标语是：早上吃好，晚上早睡。美的的标语是：大战九十天，创造新辉煌。回来后又问了一些工人当时的工资，格力大概2200元，美的只有1600多。知道这些以后，我就决定选格力而不是美的了。}

“早上吃好，晚上早睡”，呵呵，这就是企业文化了。你能从这点看到企业的差别，说明你已经是高手了。记得大概10年前，有个招行的朋友就和我谈起格力空调，说觉得有很多地方和我们公司挺像的，呵呵，我觉得他们很多地方做得比我们好，当然我们也做得还行。（2010-06-02）

\BookQuestion{网友：担心董明珠退休后接班人的问题。}

这正是我担心的地方，不然现在这价就不贵哦。

\BookQuestion{网友：董明珠很少公开谈到接班人的问题，只说过公司有很多自己培养的年轻人很优秀，如果她打算5年后退休，现在应该还在考察阶段吧。段总以前提到一个逻辑，好的企业文化——好的董事会——好的接班人……这种机制要维护文化不容易。（2010-07-07）}

\BookQuestion{网友：我认为格力没有好的文化传承。如果董明珠退休了，格力还能有这份坚持吗？}

那是你买的时候要考虑的东西。（2014-08-21）

% 来源：EPUB part096.xhtml
\BookMinorHeading{UT斯达康}

\BookQuestion{“UT斯达康第一财季净亏损1030万美元。”（新闻）}

神奇的UT斯达康，这么久了居然还没亏完。N年前曾经糊里糊涂买过不少这只股票，然后去拜访过一次公司，见过当时的CEO，回来后就用最快的速度把股票都卖了。好像还赚了点小钱。要是留到现在可就亏大了。现在居然还有这么大的成交量，很好奇都是谁还在买呢？

\BookQuestion{网友：为什么卖？拜访过程中你发现什么了？}

呵呵，不方便细说，但至少六顶帽子里有一顶以上出来说NO了。（2011-05-07）

六顶帽子思考法，呵呵，网上很容易就找到了。（2011-04-08）

\BookQuestion{网友：您访问企业时主要看哪些方面？怎样才能看出价值呢？}

这个我也不会，至少没办法通过参观企业看出好企业来，但有时可以看出来有些企业不太好。有很多企业你参观完了就不想投了。（2010-10-05）

% 来源：EPUB part097.xhtml
\BookMinorHeading{脸书}

\BookQuestion{网友：脸书的商业模式，我个人理解是在链接人的同时，通过用户数据分析了解人，并据此通过广告等形式商业化。如果个人数据的运用权更多地交给用户选择，那么基于信息流推送的商业模式是否还可以成立？}

个人认为脸书（FB）还是一样厉害。用户没有别的选择。（2018-05-01）

\BookQuestion{网友：FB高管的不断离职会不会对公司产生致命打击？现价可以买入做长期投资吗？}

我不推荐人买股票，我觉得如果你需要问人该不该买的时候那就还没到你能买的时候。我个人觉得FB也是家不错的公司，生意模式和文化都不错，扎克伯格年轻有为，还能好好干很多年。公司有人离职很正常。（2019-03-20）

\BookQuestion{网友：您常说本分的重要性，FB近几个月频繁爆出隐私泄露的问题，您怎么看？}

FB原来的使命是“连接世界”，但没有注意隐私。现在使命不变，但注重隐私已经是前提了。好公司犯错误时往往是买入的机会。（2019-03-28）

隐私是头等事，扩张是正经事。（2019-05-01）

脸书准备卖掉了，有些不舒服的东西。（2020-07-31）

今天清空了FB，算是赚了不少。由于自己不用FB的任何产品，所以对FB的理解一直没达到很透彻的地步。最近一段时间，抖音就把它吓成这个样子，说明它可能还是有弱点的。我的卖出也符合“对自己不够了解的公司，涨了也想卖跌了也想卖”的情况。（2020-08-04）

\BookQuestion{网友：为什么大道认为看不太懂FB的商业模式呢？}

看不懂需要理由吗？我还以为看懂才需要理由呢。我个人完全不用FB以及FB的任何产品。

\BookQuestion{网友：为什么处理掉呢？}

为什么？就是觉得不太喜欢了呗。看到扎克伯格在国会作证时说的话后，我就“粉转路”了。（2022-02-03）

FB真的是我糊里糊涂买、糊里糊涂赚，但很清楚卖掉的公司之一吧。我还记得当时看到“小扎”在国会质询回答问题时，我突然感觉到FB的护城河其实很浅，面对抖音他们完全没有招架能力，连阴招都上了。看完第二天就基本清仓了，只有几万股是留到年底卖掉的。话说回来，FB本来买得也不太多，所有管的账号大概能有2\%，卖光了也没啥不开心的。（2022-10-09）

\BookQuestion{网友：因为企业文化不好？}

是的，我有点不太喜欢他们的企业文化。

\BookQuestion{网友：我查到了这一段，时间也跟大道表态准备卖掉FB的时间相吻合——“到了2020年扎克伯格的假面具摔得更彻底，七月时扎克伯格在美国国会举行听证会向中国企业‘泼脏水’，声称‘中国从美国科技公司’窃取技术‘证据确凿’。而对于这一问题，苹果总裁、谷歌总裁、亚马逊总裁的回答均为‘没有发生过此类事件’。”就是这个事情后就决定卖了。（2022-02-10）}

% 来源：EPUB part098.xhtml
\BookMinorHeading{双汇}

\BookQuestion{网友：如何避免投资碰到黑天鹅事件？}

黑天鹅事件在买的时候就应该避开了，不然你就不应该投资。

价值投资者其实最喜欢看到黑天鹅事件的发生，这种时候往往会有机会。上一次金融危机无数人中招，但巴菲特那么大盘子，那些一掉到底不复回的股票他一个都没有就是一个例子，这可绝对不是运气哈。顺便说一句，那些股票我也一个都没有。

要是你的“估值”会和公司每一个产品的推出和定价联系在一起，你早晚会在不赚钱的前提下累趴下的。

补充一点：买的时候避开的意思就是不懂不做，事前想清楚。

去年有个朋友想投双汇，问我的看法。我说我不懂，你好像也不比我懂，但我觉得中国食品行业风险大，不知道啥时候就会出点啥问题，结果不幸而言中了。（2011-03-30）

\BookQuestion{网友：目前双汇的“瘦肉精”事件，到底如何判断目前的危机是危险还是机会呢？}

这种时候关键看你对公司及其文化有多了解。比如，GE的杰夫·伊梅尔特出来说公司没问题时，我是相信他确实认为自己公司没问题，所以我就敢下手买。

如果王石碰到类似问题时出来说话我也会相信的。

但如果是雷曼或美国国际集团（American International Group，AIG）的CEO当时出来说，我是无法相信的，因为不了解。

你如果了解双汇并绝对相信他们的话，那就会是个机会，不然还是离远一点好。

阿里的供应商欺诈问题反映了当时管理层的心态已经出问题了，但马云后来的处理还是很给力的，如果股价因此大掉的话，我觉得会是个机会。不过，实际情况好像掉得也很少，看不出有大油水。

我个人认为，一般情况下，急跌的股票还是不要碰的好。资本市场大多数情况下都很聪明，但偶尔的错判往往会是机会。

投资不要怕错失机会，最重要的是不要犯大错。没听说过这种事有规律可循的。（2011-03-24）

\BookQuestion{网友：当企业出现什么样的问题时您会卖掉，或继续观察或买入？}

你只要能分清做了错的事情和把事情做错了，你就会有结论的。

(2011-03-17)

\BookQuestion{网友：请举个例子。}

我个人的经验基本是用“己所不欲”排除法来鉴别企业的，比如一家明知故犯地加三聚氰胺的公司我就会认为是一家不好的公司。又比如一家为500大而500大的公司我就会很小心。（2010-03-21）

% 来源：EPUB part099.xhtml
\BookMinorHeading{比亚迪}

\BookQuestion{网友：芒格为什么买比亚迪？}

我现在觉得有点明白芒格为什么投比亚迪了。如果我能有当时一样的条件，我也可以投。

大概就是个未来现金流折现的问题。芒格投20亿（的1/10）在40岁出头的王传福身上，认为他未来20年（或更长）肯定可以给出多倍回报。这么看王传福的话，当时那个价确实便宜。（2010-06-21）

\BookQuestion{网友：不明白芒格为何如此看好比亚迪。}

我也一直认为买比亚迪不是巴菲特的风格。无论如何，芒格的面子是要给的，结果也很好。（2010-09-29）

\BookQuestion{网友：比亚迪回归A股，是A股价值投资者的一次机会吗？}

你问的是什么机会？

到目前为止我还是没完全搞懂比亚迪。我觉得王传福是个对商业机会很敏感且非常用功的人，但不是很明白为什么他们现在还在四面出击，同时还要借那么多钱。

这次在巴菲特那里还和王传福一起吃了个饭并聊了一会，但时间有限，没能了解太深，有待进一步了解。如果比亚迪的储能柜真的能做成的话，那可是大事，因为目前利用太阳能最大的问题就是还不能有效储存能量。可是，比亚迪现在的业务很散，给我的感觉好像是没足够信心集中精力在能源产品上的样子。

另外，我不太喜欢他们的中国第一和世界第一的那些个目标，那些个东西消费者其实不关心。你看看苹果什么时候说过要提高市场占有率了？苹果总是说且只说要做最好的产品（消费者最喜欢的产品）。

为了达到第一的目标（实际上是数量上的第一），企业一定会做些很奇怪的事情，比如价格战，比如因为过于急于上量而导致的质量问题等。当我看到他们说把2015年成为第一的目标推迟到2018年的时候，我就觉得他们还没有明白过去在哪里出了问题。当然，如果王传福真是像芒格说的那样，是个爱迪生和韦尔奇的合体超人（超人是我加的），那可能一切就都是对的。但这些超出了我的理解能力范围，所以我自己不会去抓这个机会。（2011-05-11）

\BookQuestion{网友：对巴菲特买比亚迪是买王传福未来几十年现金流能理解一些了，以前单纯看比亚迪年报看不出什么，结合着比亚迪几十年的发展史来看，那个价格买王传福优秀卓越的经营能力确实不贵。}

呵呵，我想巴菲特买的也是王传福的未来现金流（的折现）。（2011-01-09）

% 来源：EPUB part100.xhtml
\BookMinorHeading{万科}

\BookQuestion{网友：您和王石接触过吗？}

没有很深入打交道的机会，有过几次聊天或吃饭的机会。对他印象还不错，有点“青涩”的味道。（2010-05-17）

那个对王石的评语不是我说的，是他自己说的。不过我觉得“青涩”很难得，现代人普遍太圆滑了。（2010-06-06）

\BookQuestion{网友：有家公司，总市值263亿，净资产160亿，净利润12亿，这只是商业地产股，对你有吸引力吗？}

如果你能保证里面没假账，净资产不贬值，利润不下降还不会用途不当，这个价格倒也不贵。呵呵，了解一家公司不容易就是在这些地方。当然，有时了解也许没那么复杂。

我记得当年我们买万科时就有人问过我，说万一万科假账怎么办。

我说，以我认识的王石而言，他绝不是会关在房间里和财务商量个假账来蒙股东的人。其实那时和王石不熟，现在也不算熟，就是直觉而已。（2010-03-14）

\BookQuestion{网友：刚听了《财经会客厅》王石的访谈。王石说万科的目标定位是健康丰盛，不是最大、第一，讲到了目标和结果，从急着登顶到坦然的体验过程，讲到了社会浮躁心态——恨不得把三十年的、几代人的事只争朝夕地干完，怎么都听过啊？和段总说过的健康长久、敢为天下后、未必最好肯定能到罗马、快即慢、平常心等等，不是一个理吗？}

能够长久经营企业的人对企业的体会实际上都差不多，不信你把马云的东西拿来比较下，你会发现有很多其实也挺像的。不过，“不幸的家庭则各有各的不幸”。（2011-01-27）
